% An explicit conditional bridge from the current reference5 surrogate
% to nearest retrieval utility. No calibration or empirical claim is added.
\section{Vì sao cần ước lượng trước khi chọn Q?}
\label{sec:formal-belief-utility-bridge}

Ước lượng ở đây không nhằm đoán lại GPS để gửi lên máy chủ. Nó phân bổ
trọng số cho những vùng có thể chứa người dùng, rồi chọn Q có khả năng
thu hồi POI của các vùng đó. Phần này chứng minh ý nghĩa chính xác của
điểm độ phủ và chỉ ra điều kiện còn thiếu để coi nó là utility kỳ vọng
thật. Không giả sử bộ lọc hiện tại đã được hiệu chuẩn.

\subsection{Điểm độ phủ là kỳ vọng trong mô hình lưới}

Giữ một sự kiện, lịch sử đã bảo vệ, catalogue và các miền chọn Q cố định.
Gọi $\ell$ là một đại diện lưới công khai, $b_\ell\geq0$ là belief và
$\sum_\ell b_\ell=1$. Dịch vụ ánh xạ tọa độ đại diện tới trạng thái truy
cập $a(\ell)$ theo cùng quy tắc công khai dùng khi nhận Q. Nếu hai làn
trùng tọa độ, $a(\ell)$ có thể khác trạng thái làn của đại diện; không
được thay âm thầm bằng khoảng cách từ làn đó.

Gọi $\mathcal C_\ell$ là các category có reference nearest không rỗng
tại $a(\ell)$ và $R_{\ell,c}$ là reference top5 của category $c$.
Mỗi danh sách gồm các ID khác nhau và có thể ít hơn5 POI. Trọng số mà
mô hình gốc tạo là
\begin{equation}
 a_{\ell p}=\begin{cases}
 \displaystyle\sum_{c\in\mathcal C_\ell}
 \frac{\mathbf1\{p\in R_{\ell,c}\}}
 {|\mathcal C_\ell|\,|R_{\ell,c}|},&|\mathcal C_\ell|>0,\\[2mm]
 0,&|\mathcal C_\ell|=0,
 \end{cases}
 \qquad w_b(p)=\sum_\ell b_\ell a_{\ell p}.
 \label{eq:formal-belief-reference-weights}
\end{equation}
Một hàng có reference có tổng trọng số1; hàng rỗng có tổng0. Như vậy
$\sum_p w_b(p)$ có thể nhỏ hơn1 khi belief đặt khối lượng vào vùng
không có reference. Đây không phải lỗi chuẩn hóa cần bù bằng POI giả.

Gọi $A_{10}(Q)$ là hợp signature top10 của mọi category tại các Q,
được lấy theo cùng phép ánh xạ tọa độ của dịch vụ. Đặt
\begin{equation}
 f_Q(\ell)=\sum_p a_{\ell p}\mathbf1\{p\in A_{10}(Q)\},
 \qquad F_b(Q)=\sum_\ell b_\ell f_Q(\ell).
 \label{eq:formal-belief-expected-cover}
\end{equation}
Đổi thứ tự hai tổng hữu hạn cho
$F_b(Q)=\sum_p w_b(p)\mathbf1\{p\in A_{10}(Q)\}$, đúng mục tiêu ở
(\ref{eq:formal-service-cover}). Khi $\mathcal C_\ell$ không rỗng,
\begin{equation}
 f_Q(\ell)=\frac{1}{|\mathcal C_\ell|}
 \sum_{c\in\mathcal C_\ell}
 \frac{|R_{\ell,c}\cap A_{10}(Q)|}{|R_{\ell,c}|},
 \qquad 0\leq f_Q(\ell)\leq1.
 \label{eq:formal-belief-category-recall}
\end{equation}
Do đó $F_b$ là kỳ vọng của \emph{độ thu hồi reference nearest trung bình
đều giữa các category không rỗng}, theo mô hình lưới $b$. Nó không dùng
trực tiếp category hoặc purpose riêng của người dùng.

Với dịch vụ tĩnh, cùng catalogue, thứ tự tie và phép ánh xạ tọa độ,
phản hồi $L=30$ chứa prefix top10. Nếu local ranking được đánh giá đúng
tại $a(\ell)$, Tính chất~3 ở Mục~\ref{sec:formal-service-local-exactness}
cho Recall nearest trung bình trên các category đó không thấp hơn
$f_Q(\ell)$. Khẳng định này chỉ nối surrogate với \emph{reference tại
đại diện lưới}. Nó chưa nối tới GPS liên tục, làn thật khác trạng thái
truy cập, fastest, radius hay detour; và không đổi mẫu số khi reference
rỗng. Trong phương trình, $f_Q(\ell)=0$ ở hàng rỗng chỉ biểu thị đóng góp
surrogate bằng0, không quy ước Recall benchmark của reference rỗng bằng0.
\footnote{Công thức trọng số: \path{benchmark/anchor_belief.py},
\texttt{PublicAnchorModel.\_poi\_weights} và
\texttt{AnchorBelief.poi\_weights}. Wrapper
\path{benchmark/response_aware_belief.py} giữ reference5 gốc, dùng
signature planner10 và kiểm tra catalogue/phép truy cập phù hợp.}

\subsection{Sai lệch belief ảnh hưởng cận chọn Q như thế nào?}

Giả sử có một phân bố mục tiêu $\mu$ trên \emph{cùng lưới và cùng
reference/category weighting}. Giữ mọi miền khả thi như trên và giả
sử sai lệch total variation có một cận đã được xác lập riêng:
\begin{equation}
 \operatorname{TV}(b,\mu)=\frac12\sum_\ell|b_\ell-\mu_\ell|
 \leq\zeta.
 \label{eq:formal-belief-calibration-assumption}
\end{equation}
Đây là \textbf{giả định có điều kiện}, không phải giá trị đã đo của bộ
lọc hiện tại. Vì mọi $f_Q$ thuộc $[0,1]$, ta có đồng thời cho mọi Q:
\begin{equation}
 |F_b(Q)-F_\mu(Q)|\leq\zeta.
 \label{eq:formal-belief-expectation-error}
\end{equation}
Để thấy điều này, tổng các phần dương của $b-\mu$ và trị tuyệt đối tổng
các phần âm đều bằng TV. Nhân từng phần với một số trong $[0,1]$ giữ
hiệu hai tổng trong đoạn $[-\operatorname{TV},\operatorname{TV}]$.

Đặt $\sigma=\delta_F$ là cùng \texttt{utility\_slack} ở
(\ref{eq:formal-service-slack-floor}), hiện bằng0,03 điểm surrogate.

\begin{samepage}
\noindent\textbf{Hệ quả (cận độ phủ khi belief có sai lệch).}
Nếu greedy và slack lý tưởng thỏa
$F_b(Q_{\rm alg})\geq\tfrac12\max_Q F_b(Q)-\sigma$,
thì
\begin{equation}
 F_\mu(Q_{\rm alg})\geq
 \max\!\left\{0,\ \tfrac12\max_Q F_\mu(Q)-\sigma-
 \tfrac32\zeta\right\}.
 \label{eq:formal-belief-utility-bridge}
\end{equation}
\end{samepage}

\noindent\emph{Chứng minh.}
Gọi $Q^*_\mu$ là tối ưu cho $\mu$ trên cùng tập khả thi. Dùng
(\ref{eq:formal-belief-expectation-error}) trước và sau cận greedy:
\[
 \begin{aligned}
 F_\mu(Q_{\rm alg})
 &\geq F_b(Q_{\rm alg})-\zeta\\
 &\geq\tfrac12F_b(Q^*_\mu)-\sigma-\zeta\\
 &\geq\tfrac12F_\mu(Q^*_\mu)-\sigma-\tfrac32\zeta.
 \end{aligned}
\]
Độ phủ không âm nên có thể lấy thêm cận0. \hfill$\square$

Cận giải thích tại sao có layer ước lượng: giảm sai lệch của phân bố
dùng chọn Q có thể cải thiện mục tiêu dịch vụ; chỉ giảm lỗi của một tọa
độ trung bình không đủ. Tuy nhiên, $\zeta$ hiện chưa được xác lập, nên
không thay số để tuyên bố cận Recall thực nghiệm. Muốn áp dụng cho người
dùng thật còn cần kiểm tra sai lệch đại diện lưới, phân bố category và
quan hệ giữa sensor với trạng thái xếp hạng. Cận vẫn chỉ cho một sự
kiện trên miền khả thi đã cố định; sự thay đổi Q làm đổi miền ở các sự
kiện tiếp theo và không cho một bảo đảm tối ưu toàn chuyến.
